\documentclass{exam}
\usepackage{../../shah311}

\hypersetup{colorlinks=true, linktoc=section, linkcolor=blue}

% p.140: 4.5.3, 4.5.4, 4.5.6, 4.5.7, 
% p.141: 4.6.4, 4.6.5, 4.6.7.

\pagestyle{headandfoot}
\firstpageheadrule
\runningheadrule
\firstpageheader{Parikh \\ Intro Topology}{Homework 11}{Jeevan Shah}
\runningheader{Intro Topology \\ Homework 1}{}{Shah}
\firstpagefooter{}{}{}
\runningfooter{ }{\thepage}{ }

\printanswers

\begin{document}
\colorbox{red}{\underline{1:}}
\begin{solution}
\begin{proof}
    We must show that $A$ is open in $X$. Since $\forall\,x \in A\, \exists U$ open such that $x \in U$ and $U \subset A$ we know that 
    \[
        \bigcup_{i \in I} U_i = A.
    \]
    Since each $U_i \subset A \subset X, U_i \subset X$ and $A \subset \bigcup U_i$ since every $x \in A$ is also in some $U_{i}$ and thus the union. So $A$ must be open in $X$ since the topology is closed under arbitrary union.
\end{proof} 
\end{solution}

\colorbox{red}{\underline{3:}}
\begin{solution}
\begin{proof}
    \begin{parts}
        \part We must show that $\mathcal{T}_C$ is a topology. First by definition of $\mathcal{T}_C$ both $X, \varnothing \in \mathcal{T}_C$. Now, if $\braces{U_i}_{i \in I} \in \mathcal{T}_C$ to show that $\bigcup U_i \in \mathcal{T}_C$ we see 
        \[
            X - \bigcup{U_i} = \bigcap\left(X - U_i\right)
        \]
        which must be countable since each $X - U_i$ is either all of $X$ or countable (by definition of $\mathcal{T}_C$) and their intersection must also be countable as a subset of countable sets. If $U, V \in \mathcal{T}_C$ are nonempty then $X - \left(U \cap V\right) \in \mathcal{T}_C$ since $U \cap V$ is countable (by being a subset of countable sets) and so $X - \left(U \cap V\right)$ must also be countable. Thus $\mathcal{T}_C$ is a topology. 

        \part No, $\mathcal{T}_{\infty}$ is not a topology on $X$ because the intersection of infinite sets is not always infinite. Take $X = \ZZ$ and let $U = \ZZ^{+} \cup \braces{0}$ and $V = \ZZ^{-} \cup \braces{0}$. It is clear that $U, V \in \mathcal{T}_{\infty}$ but $U \cap V = \braces{0} \not\in \mathcal{T}_\infty$ and $\mathcal{T}_\infty$ is not closed under finite intersections and thus cannot be a topology.
    \end{parts}
\end{proof}    
\end{solution}

\colorbox{red}{\underline{5:}}
\begin{solution}
    \begin{proof}
        Let $\mathcal{T}_{\mathcal{A}}$ be the topology generated by $\mathcal{A}$ and let $\mathcal{T}'$ be the intersection of all topologies $\mathcal{T}$ of $X$ that contain $\mathcal{A}$. It is clear by definition that $\mathcal{T}' \subseteq \mathcal{T}_{\mathcal{A}}$ since $\mathcal{T}_\mathcal{A}$ is apart of the collection of sets being intersected.Now, consider some $U \in \mathcal{T}_{\mathcal{A}}$. Since $\mathcal{A}$ is a basis, $U = \bigcup A_i$ for $A_i$ basis elements. Now let $\mathcal{T}$ be any topology on $X$ containing $\mathcal{A}$ Since $A \subseteq \mathcal{T}, \mathcal{A}_i \subseteq \mathcal{T}$ and $\mathcal{T}$ is closed under arbitrary union so $U \in \mathcal{T}$ thus $\mathcal{T}_{\mathcal{A}} \subseteq \mathcal{T}'$ and $\mathcal{T}_{\mathcal{A}} = \mathcal{T}'$. 

        Now, let $\mathcal{A}$ be a subbasis and $\mathcal{T}_\mathcal{A}$ be the topology generated by $\mathcal{A}$. As before it is clear that $\mathcal{T}' \subset \mathcal{T}_\mathcal{A}$ since $\mathcal{T}_\mathcal{A}$ belongs to the family of sets being intersected. We define $\mathcal{B}$ to be the collection of all finite intersection of elements of $\mathcal{A}$. Then $\mathcal{T}_\mathcal{A} = \mathcal{T}_\mathcal{B}$ where $\mathcal{T}_\mathcal{B}$ is the topology generated by $\mathcal{B}$. It follows that $\mathcal{B} \subseteq \mathcal{T}$ since $B = \bigcap A_i$ and each $A_i \subset A \in \mathcal{T}$. Applying the idea proven above, $\mathcal{T}_\mathcal{B} = \mathcal{T}_\mathcal{A} = \mathcal{T}'$ since $\mathcal{T}_\mathcal{B}$ is generated by a basis.
    \end{proof}
\end{solution}

\colorbox{red}{\underline{7:}}
\begin{solution}
    We know that $\mathcal{T}_1 \subseteq \mathcal{T}_2$ since any basis element $(a, b) \in \mathcal{T}_1$ is automatically a basis element in $\mathcal{T}_2$ by definition. $\mathcal{T}_1 \subseteq \mathcal{T}_4$ as well since for $(a, b) \in \mathcal{T}_1$ and $x \in (a, b)$, there exists $x \in (a, x] \subset (a, b)$ with $(a, x] \in \mathcal{T}_4$. $\mathcal{T}_3 \subseteq \mathcal{T}_1$ since for any $U \in \mathcal{T}_3$ finite, there exists some $\varepsilon$ such that $(y - \varepsilon, y + \varepsilon) \subseteq \RR \setminus F = U$ for $F = \RR \setminus U$ where $(y - \varepsilon, y + \varepsilon) \cap F = \varnothing$. It is clear that $\mathcal{T}_5 \subseteq \mathcal{T}_1$ since for $x \in (a, b) \in \mathcal{T}_1$, $x \in (-\infty, b) \in \mathcal{T}_5$ and $(a, b) \subseteq (-\infty, a)$. It follows that $\mathcal{T}_5 \subseteq \mathcal{T}_2$ and $\mathcal{T}_5 \subseteq \mathcal{T}_4$ from transitivity of subsets and $\mathcal{T}_3 \not\subseteq \mathcal{T}_5$ and $\mathcal{T}_5 \not\subseteq \mathcal{T}_3$ so $\mathcal{T}_3$ and $\mathcal{T}_5$ are not comparable since the complement of $\left(-\infty, a\right)$ is always infinite and the union of infinite sets will always be infinite. $\mathcal{T}_2 \subseteq \mathcal{T}_4$ since for $x \in (a, b) \in \mathcal{T}_2, x \in (x - \varepsilon, x + \varepsilon] \subseteq (a, b)$ for epsilon small so $(a, b) \in \mathcal{T}_4$. For $x \in (a, b) \setminus K$ find $\max\braces{K \cap (0, x)} < c < x$ so that $x \in (c, x] \subseteq (a, b) \setminus K$. Now, applying associativity of subsets we can see $\mathcal{T}_3 \subseteq \mathcal{T}_2$ and $\mathcal{T}_3 \subseteq \mathcal{T}_4$.
\end{solution}

\colorbox{red}{\underline{8:}}
\begin{solution}
    \begin{parts}
        \part We must show $\mathcal{C}$ is a collection of open sets such that for each $U \subset \RR$ each $x \in U\, \exists C \in \mathcal{C}$ such that $x \in C \subset U$. Consider some $U \subset \RR$ open and $x \in U$. Then take $C = (x - \varepsilon, x + \varepsilon)$ for $\varepsilon$ small and $x \pm \varepsilon \in \QQ$ so $C \in \mathcal{C}$ and then $C \subset U$. Thus, by Lemma 13.2 $\mathcal{C}$ is a basis for $X$. By Lemma $13.1$ $\mathcal{T}$ is the union of all elements of $\mathcal{B}$ which is exactly the standard topology. 

        \part It suffices to show that $\mathcal{T}_\mathcal{C}$ (the topology generated by $\mathcal{C}$) is strictly coarser than $\mathcal{T}_\ell$, i.e., $\mathcal{T}_\mathcal{C} \subset \mathcal{T}_\ell$. Consider some basis element $[a, b)$ of $\mathcal{T}_\mathcal{C}$ and $x \in [a, b) \in \mathcal{T}_\ell$ since $[a, b)$ is also a basis element of $\mathcal{T}_\ell$. It follows that $\mathcal{T}_\mathcal{C} \subset \mathcal{T}_\ell$. We must now show $\mathcal{T}_\ell \not\subseteq \mathcal{T}_\mathcal{C}$. Consider $[c, d) \in \mathcal{T}_\ell$ a basis element with $c \in \RR \setminus \QQ$ and let $x = c \in [c, d)$. Then there is no open set in $\mathcal{T}_\mathcal{C}$ such that $x$ is contained in it and it is a subset of $[c, d)$ since all endpoing of $\mathcal{C}$-basis elements are rational. Thus $\mathcal{T}_\ell \not\subseteq \mathcal{T}_\mathcal{C}$ and $\mathcal{T}_\ell$ is strictly finer than $\mathcal{T}_\mathcal{C}$.
    \end{parts}
\end{solution}
\end{document}